\scanpage{147}
\mspage{475}
% position: upper paragraph
Si quelque chose ici sort du commun, ce n'est pas bien sûr la simple
\mbox{p r é s e n c e} de ce double mécanisme, ni non plus (il me semble) la force de
la composante ``répressive'' à proprement parler, la force donc de la répression
des traits, attitudes, pulsions ``yin''. Il n'y a pas de commune mesure ici avec
ce qui avait lieu chez ma mère, dont la vie (et celle de ses proches) a été dé-\linebreak
vastée par sa haine (restée occulte sa vie durant) de ce qui faisait d'elle une
femme. A aucun moment, je crois, mes façons d'être n'ont été entièrement \unsure{exempte|exemptes}
d'une certaine douceur, voire de tendresse, qui obstinément arrondissaient les
angles du personnage que je m'étais taillé depuis mon enfance, et qui \unsure{m'attiraie|m'attiraient}
souvent sympathie et affection. Le côté exceptionnel se trouverait plutôt dans \unsure{l|la}
\mbox{d é m e s u r e} de mes investissements, dans la démesure de l'énergie que j'in\linebreak
vestis dans mes tâches, sans m'en laisser distraire par un regard à droite ni à
gauche! En dehors du travail proprement dit, mon esprit continuellement est pro-\linebreak
jeté vers l'accomplissement, vers l'aboutissement de telle ou telle étape du tra\linebreak
vail. Cette attitude-là (``Zielgerichtetheit'' en allemand, ``aimdirectedness'' en
anglais) est par excellence une attitude yang, une attitude de \mbox{t e n s i o n},
de \mbox{f e r m e t u r e} à tout ce qui n'apparaît pas directement lié à la tâche

% position: middle paragraph
Cette démesure était de nature à susciter en autrui l'image d'une sorte de
``super-man'' ou ``super-mâle'', admirable certes hélas! (vu les valeurs qui ont
cours), mais suscitant aussitôt (à un niveau qui reste inconscient le plus sou-\linebreak
vent) d'instinctives réactions de défense voire d'antagonisme devant un tel dé-\linebreak
ployement de force, ressenti comme menaçant voire agressif, ou en tous cas dan-\linebreak
gereux \(({}^{108}2)\). Et surtout, cette image irrésistiblement évoque l'image du
``\mbox{s u p e r - p è r e}'', et met en branle aussitôt la multiplicité ambiguë des
réactions d'attraction et de répulsion nouées autour du sempiternel conflit au
père... C'est là \mbox{m a} contribution dans ces relations d' \mbox{a m b i g u i t é},
qui ont été si communes dans ma vie, et auxquelles je me suis retrouvé confronté
tant de fois au cours de Récoltes et Semailles. Cette ambiguïté est renforcée,
non diminuée, par la persistance de traits yin en moi qui alimentent une sympa-\linebreak
thie, que la seule hypertrophie des traits yang en une sorte de gigantesque
``superman'' serait impuissante à susciter.

% position: lower paragraph
Et à nouveau je peux constater, dans ces sempiternelles ``relations d'ambi-\linebreak
guité'', que je ne fais encore que récolter ce que j'ai moi-même semé, même si à
chaque fois la récolte s'avère inattendue (et ma\inserted{l}venue...)! Car la motivation (\unsure{o|ou}
du moins \mbox{u n e} des motivations) qui pousse ``le patron'' en moi à se dépasser
sans cesse dans l'accumulation des oeuvres, n'a-t-elle pas été justement de \unsure{forc|forcer}
et relancer sans cesse l'estime de mes pairs (en premier lieu) et de mes impairs
% position: foot
\margin{146}
